\chapter{Closed Systems and Memory: Finite Populations, MVA, Campbell's Theorem, Eviction}\label{sec:L4}

All models so far were \term{open systems}. Customers arrive from an infinite
outside population, leave after service, and the arrival rate does not change
when the queue is long. Agent sessions are different. A session submits a turn,
receives the response, runs a tool, and only then submits its next turn. A
waiting session cannot create a new turn, so the longer the queue, the fewer the
arrivals. This chapter treats this \emph{closed} structure
(\cref{L4:sec:closed}), computes the mean memory held by sessions
(\cref{L4:sec:campbell}), and solves the combinatorial optimisation problem of
what to drop when memory runs short (\cref{L4:sec:knapsack}).

\section{Closed systems and mean value analysis}\label{L4:sec:closed}

The simplest closed system (\cref{L1:sec:systems}) is the finite-source model
M/M/1//$N$ (\cref{L1:def:queues}(f)). $N$ customers (sessions) alternate between a
delay station, where each thinks for an exponential time of mean $Z$, and a FIFO
server with service times $\mathrm{Exp}(\mu)$: a session submits a turn when it
finishes thinking and starts thinking again when the turn is served.

\serving{Sessions are the customers of the closed loop. The think time is the
time for tool execution and user response, and the server is the prefill stage.
A session waiting for prefill cannot submit its next turn. A KV memory pool holds
a few million tokens and real agent contexts reach hundreds of thousands of
tokens, so a replica holds about 10 live sessions at a time. At this scale the
$N=\infty$ approximation can be badly wrong.}

Let $X(t)$ be the number of turns at the server (waiting or in service) at time
$t$. There are $N-X(t)$ thinking sessions, each finishing its think time
independently at rate $1/Z$ (\cref{prop:memoryless}), so $X$ is a birth--death
process on $\{0,1,\dots,N\}$ with birth rates $\lambda_n=(N-n)/Z$ and death
rates $\mu_n=\mu$ ($n\ge1$). By \cref{thm:birthdeath} the stationary law is
\begin{equation}\label{L4:eq:piN}
  \pi_N(n)=\frac{1}{G_N}\prod_{k=0}^{n-1}\frac{(N-k)/Z}{\mu}
  =\frac{1}{G_N}\,\frac{N!}{(N-n)!}\Bigl(\frac{1}{\mu Z}\Bigr)^{n},
  \qquad n=0,1,\dots,N,
\end{equation}
with normalising constant $G_N$. The state space is finite, so no stability
condition is needed.

In an open system PASTA (\cref{thm:pasta}) told us that an arriving customer sees
the time-average law. In a closed system arrivals are not Poisson, and an
arriving customer does not see the time-average law. Instead the following
holds.

\begin{theorem}[Arrival theorem \cite{reiser1980,lazowska1984}]\label{thm:arrival}
In M/M/1//$N$ in steady state, the law of the number of turns at the server seen
by an arriving turn (excluding itself), that is, the long-run fraction $a_N(n)$
of arrivals that occur in state $n$, equals the time-stationary law of the same
system with $N-1$ sessions: $a_N(n)=\pi_{N-1}(n)$, $n=0,\dots,N-1$.
\end{theorem}

\begin{proof}
In state $n$ arrivals occur at rate $(N-n)/Z$. By the ergodic theorem
(\cref{L2:thm:ergodic}) and the law of large numbers for jump counts, the number of
arrivals in state $n$ during $[0,t]$, divided by $t$, converges almost surely to
$\pi_N(n)(N-n)/Z$. Hence
\[
  a_N(n)=\frac{(N-n)\,\pi_N(n)}{\sum_{m=0}^{N}(N-m)\,\pi_N(m)} .
\]
By \eqref{L4:eq:piN}, for $n\le N-1$ we have
$(N-n)\frac{N!}{(N-n)!}=N\frac{(N-1)!}{(N-1-n)!}$, and for $n=N$,
$(N-n)\pi_N(n)=0$. Therefore
$a_N(n)\propto\frac{(N-1)!}{(N-1-n)!}(\mu Z)^{-n}\propto\pi_{N-1}(n)$, and since
both are probability distributions they are equal.
\end{proof}

\begin{remark}[Back to PASTA as $N\to\infty$]\label{L4:rem:arrival-limit}
Fix $\lambda>0$ and $\mu>\lambda$, and let the think time grow with the population,
$Z=N/\lambda$, so that the total arrival rate when all sessions think stays $\lambda$.
Then $\pi_N$ and $\pi_{N-1}$ both converge to the M/M/1 law $(1-\rho)\rho^n$
(\cref{L4:ex:limit}), so in the limit an arriving turn sees the time-average law: this
is PASTA. For finite $N$ the arrival theorem measures how far the closed system is from
that limit.
\end{remark}

The arrival theorem lets us solve a closed system by recursion on the number of
sessions. With $n$ sessions, let $Q_n$ be the stationary mean number of turns at
the server, $R_n$ the mean response time of a turn (wait plus service), and $X_n$
the throughput (turns completed per unit time).

\begin{theorem}[Mean value analysis \cite{reiser1980}]\label{thm:mva}
In M/M/1//$n$, $Q_0=0$ and for $n\ge1$
\[
  R_n=\frac{1+Q_{n-1}}{\mu},\qquad X_n=\frac{n}{Z+R_n},\qquad Q_n=X_nR_n .
\]
Hence $Q_n=n(1+Q_{n-1})/(\mu Z+1+Q_{n-1})$, the server utilisation is
$\rho_n=X_n/\mu=n/(\mu Z+1+Q_{n-1})$, and the mean waiting time of a turn
(excluding service) is $W_n=Q_{n-1}/\mu$.
\end{theorem}

\begin{proof}
By \cref{thm:arrival} an arriving turn sees on average $Q_{n-1}$ turns ahead of
it. Under FIFO all of them must finish before its service begins. By
memorylessness (\cref{prop:memoryless}) the remaining time of the turn in service
again has mean $1/\mu$, and each waiting turn has mean $1/\mu$, so the mean wait
is $Q_{n-1}/\mu$, and the response time is this plus the turn's own service
$1/\mu$, which is $R_n$. One cycle of a session takes $Z+R_n$ on average.
Applying Little's law (\cref{thm:little}) to the whole system (all $n$ sessions,
always $n$ customers) gives $n=X_n(Z+R_n)$, and applying it to the server gives
$Q_n=X_nR_n$. The utilisation is $\rho_n=X_n/\mu$ by \cref{cor:busy}, and
substitution gives the remaining formulas.
\end{proof}

\begin{proposition}[Finite population]\label{prop:finite}
Let $\gamma=\mu Z>0$ and define $Q_n$ by the recursion of \cref{thm:mva}.
\begin{enumerate}[label=(\roman*)]
\item $Q_n$ is nondecreasing in $n$, $0\le Q_n\le n$, and $Q_n$ is nonincreasing
  in $\gamma$. Moreover $\rho_N<1$ for all $N\ge1$.
\item $W_N\le\rho_N/(\mu(1-\rho_N))$. The right side is the open M/M/1 waiting
  time at the same utilisation (\cref{prop:mm1}).
\end{enumerate}
\end{proposition}

\begin{proof}
Put $g_\gamma(q)=(1+q)/(\gamma+1+q)$, so that $Q_n=n\,g_\gamma(Q_{n-1})$. On $q\ge0$, $g_\gamma$ is
nondecreasing in $q$, nonincreasing in $\gamma$, and $0<g_\gamma\le1$.

(i) $0\le Q_n\le n$ follows by induction from $g_\gamma\le 1$. $Q_1=g_\gamma(0)\ge0=Q_0$,
and if $Q_n\ge Q_{n-1}$ then $Q_{n+1}=(n+1)g_\gamma(Q_n)\ge n\,g_\gamma(Q_{n-1})=Q_n$.
For $\gamma\le \gamma'$, assuming $Q_{n-1}(\gamma')\le Q_{n-1}(\gamma)$,
$Q_n(\gamma')=n\,g_{\gamma'}(Q_{n-1}(\gamma'))\le n\,g_{\gamma'}(Q_{n-1}(\gamma))\le n\,g_\gamma(Q_{n-1}(\gamma))=Q_n(\gamma)$.
$\rho_N<1$ is equivalent to $Q_{N-1}>N-1-\gamma$, so it suffices to show $Q_n>n-\gamma$ for
all $n\ge0$ by induction. For $n=0$ this is $0>-\gamma$, and a computation gives
\[
  Q_{n+1}-(n+1-\gamma)=\frac{\gamma\,(Q_n+\gamma-n)}{\gamma+1+Q_n}>0 .
\]

(ii) Since $\rho_N<1$, the inequality is equivalent to
$Q_{N-1}(1-\rho_N)\le\rho_N$, that is, $Q_{N-1}\le\rho_N(1+Q_{N-1})=Q_N$, which is
(i).
\end{proof}

\begin{example}\label{L4:ex:two}
With $N=2$ and $Z=2/\mu$ ($\gamma=2$), $Q_1=1/3$, $\rho_2=2/(3+1/3)=3/5$ and
$W_2=1/(3\mu)$. The open formula at the same utilisation gives
$\rho_2/(\mu(1-\rho_2))=3/(2\mu)$, which is $4.5$ times the actual wait.
\end{example}

\begin{remark}
For exponential work and at the same utilisation, the open formula (\cref{thm:pk})
overestimates the wait (\cref{prop:finite}(ii)). For general work, and for the price
of a miss built on the open formula (\cref{thm:price}), this comparison has not been
proved. Moreover, at most $N$ turns exist, so $0\le L_P\le N$, and no
policy change can raise the mean number of turns at the server by more than
$N-L_P$. This cap does not depend on the work law.
\end{remark}

\begin{remark}
We assumed exponential think times, but the results depend only on their mean
$Z$. The delay station is insensitive like the PS station (\cref{thm:insens}), and
by the BCMP theorem a network of delay stations and PS stations has a
product-form stationary law \cite{bcmp1975}. The service times of a FIFO server
do not generalise in this way.
\end{remark}

\section{M/G/\texorpdfstring{$\infty$}{infinity} and Campbell's theorem}\label{L4:sec:campbell}

We now count memory. Live sessions do not wait for each other and each keeps its
own KV cache, so from the point of view of memory the sessions are customers of a
delay station.

The model is the M/G/$\infty$ system (\cref{L1:def:queues}(e)): a delay station fed
by Poisson arrivals of rate $\lambda$, with i.i.d.\ service times $S$ independent of
the arrivals.

The following lemma is a standard property of the Poisson process.

\begin{lemma}[Conditional uniformity \cite[\S2.3]{ross1996}]\label{L4:lem:uniform}
For a Poisson process of rate $\lambda$, the number of arrivals in an interval
$(a,b]$ is Poisson$(\lambda(b-a))$, and conditional on this number being $n$, the
arrival times (as an unordered set) have the law of $n$ i.i.d.\ samples from the
uniform distribution on $(a,b]$.
\end{lemma}

\begin{proposition}[M/G/$\infty$ occupancy \cite{ross1996}]\label{L4:prop:mginf}
If $\E[S]<\infty$, then in an M/G/$\infty$ system started empty at time $-t$, the
number of customers at time $0$ is
Poisson$\bigl(\lambda\int_0^t\Prob(S>u)\dd u\bigr)$, and as $t\to\infty$ it
converges in distribution to Poisson$(\lambda\E[S])$.
\end{proposition}

\begin{proof}
Let $A$ be the number of arrivals in $(-t,0]$. Conditional on $A=n$, the arrival
times are i.i.d.\ uniform (\cref{L4:lem:uniform}) and the service times are
independent of them. A customer who arrived at time $-u$ is still present at time
$0$ only if $S>u$, so each customer remains independently with probability
$p_t=\frac1t\int_0^t\Prob(S>u)\dd u$. Hence the number $X$ remaining satisfies
\[
  \Prob(X=k)=\sum_{n\ge k}e^{-\lambda t}\frac{(\lambda t)^n}{n!}\binom nk p_t^k(1-p_t)^{n-k}
  =e^{-\lambda tp_t}\frac{(\lambda tp_t)^k}{k!} .
\]
Finally $\lambda tp_t=\lambda\int_0^t\Prob(S>u)\dd u\to\lambda\int_0^\infty\Prob(S>u)\dd u=\lambda\E[S]$.
\end{proof}

That the limit depends on the law of $S$ only through its mean is the
insensitivity of the delay station. When the quantity each customer holds (for
example, KV bytes) changes over time, the next theorem gives its mean.

\begin{theorem}[Campbell's mean formula \cite{kingman1993}]\label{thm:campbell}
Let $\{T_k\}$ be the points of a Poisson process of rate $\lambda$ on $\R$, and
attach to each point an i.i.d.\ mark $Y_k\sim Y$ (with values in an arbitrary
measurable space), independent of the points. For measurable $f\ge0$,
\[
  \E\Bigl[\sum_k f(T_k,Y_k)\Bigr]=\lambda\int_{\R}\E[f(t,Y)]\dd t .
\]
\end{theorem}

\begin{proof}
First consider the sum over the points in a bounded interval $I=(a,b]$ only.
Conditional on the number of points being $n$, the pairs $(T_k,Y_k)$ are $n$
i.i.d.\ copies of $(U,Y)$ with $U\sim$ Uniform$(I)$ and $Y$ independent
(\cref{L4:lem:uniform}). Hence
\[
  \E\Bigl[\sum_{T_k\in I}f(T_k,Y_k)\Bigr]
  =\sum_{n\ge0}\Prob(A_I=n)\,n\,\E[f(U,Y)]
  =\lambda(b-a)\cdot\frac{1}{b-a}\int_a^b\E[f(t,Y)]\dd t .
\]
(The second equality uses Tonelli's theorem.) For $I_m=(-m,m]$ both sides are
nondecreasing in $m$, and the monotone convergence theorem gives the claim as
$m\to\infty$.
\end{proof}

\begin{corollary}[Mean resident KV]\label{cor:resident-kv}
Let sessions arrive in a Poisson process of rate $\Lambda$, and let session $k$
have an i.i.d.\ lifetime $T_s^{(k)}$ and a context of $K_k(u)\ge0$ tokens at age $u$ (the
lifetime and the trajectory may be dependent). In steady state (started at time
$-\infty$), the mean number of resident KV tokens at a given instant is
\[
  \Lambda\,\E\Bigl[\int_0^{T_s}K(u)\dd u\Bigr].
\]
In particular, if $K(u)\equiv\bar K$ is constant for each session, this equals
$\Lambda\bigl(\E[T_s]\E[\bar K]+\Cov(T_s,\bar K)\bigr)$. It exceeds the naive
product $\Lambda\E[T_s]\E[\bar K]$ if and only if $\Cov(T_s,\bar K)>0$, that is,
when long sessions have large contexts.
\end{corollary}

\begin{proof}
Attach the mark $Y_k=(T_s^{(k)},K_k(\cdot))$ and put
$f(t,(T,K))=K(-t)\1\{0\le -t<T\}$. The resident KV at time $0$ is
$\sum_kf(T_k,Y_k)$. By \cref{thm:campbell} and Tonelli, its mean is
\[
  \Lambda\int_{-\infty}^0\E[K(-t)\1\{-t<T_s\}]\dd t=\Lambda\E\Bigl[\int_0^{T_s}K(u)\dd u\Bigr].
\]
If $K\equiv\bar K$, the integral is $T_s\bar K$, and
$\E[T_s\bar K]=\E[T_s]\E[\bar K]+\Cov(T_s,\bar K)$.
\end{proof}

Taking $K\equiv1$ gives the mean number of sessions $\Lambda\E[T_s]$, in agreement
with Little's law.

\section{Eviction: the covering knapsack problem}\label{L4:sec:knapsack}

When memory runs short, part of the KV of the paused programs must be dropped
(\term{eviction}). Program $i$ occupies $c_i>0$ bytes, and dropping it has an
expected cost $w_i\ge0$ (for example, the resume probability $p_i$ times the price
of a miss, \cref{thm:price}).

\begin{definition}[Covering knapsack problem]\label{L4:def:knapsack}
For a finite set $I$ and a target $\Delta C>0$, the \term{covering knapsack
problem} is
\[
  \min_{S\subseteq I}\ w(S):=\sum_{i\in S}w_i\quad\text{s.t.}\quad c(S):=\sum_{i\in S}c_i\ge\Delta C .
\]
Allowing $x\in[0,1]^I$ in place of $S$ and minimising $\sum w_ix_i$ subject to
$\sum c_ix_i\ge\Delta C$ gives the \term{fractional version}, which assumes that cost is
linear in the evicted fraction. We call $v_i=w_i/c_i$ the \term{price per byte}.
\end{definition}

\paragraph{Tail blocks are not fractional items.} Engines evict blocks from the tail
of a cached prefix, not whole programs. Dropping the last $d$ of $K$ cached tokens adds
\[
  r_K(d)=P(n+d,K-d)-P(n,K)=a\,d+b\bigl(Kd-d^2/2\bigr),\qquad 0\le d\le K,
\]
of prefill work. The marginal cost $a+b(K-d)$ falls as $d$ grows, so the cost is
concave in $d$, not linear. The threshold rule below is therefore exact for whole
programs and for the linear relaxation, but it does not prove that greedy block
eviction is optimal. \Cref{prop:blind} below also concerns whole programs.

\begin{theorem}[Threshold rule \cite{dantzig1957}]\label{thm:threshold}
\begin{enumerate}[label=(\roman*)]
\item For every $\theta\ge0$, the set $T_\theta=\{i:v_i\le\theta\}$ costs no more
  than any $S$ with $c(S)\ge c(T_\theta)$.
\item In the fractional version, the solution that sets $x_i=1$ in increasing
  order of $v_i$ and fills the last item partially to meet the target exactly is
  optimal (provided $c(I)\ge\Delta C$).
\end{enumerate}
\end{theorem}

\begin{proof}
(i) On $S\setminus T_\theta$ we have $w_i>\theta c_i$, and on $T_\theta\setminus S$
we have $w_i\le\theta c_i$, so
\[
  w(S)-w(T_\theta)=\sum_{S\setminus T_\theta}w_i-\sum_{T_\theta\setminus S}w_i
  \ge\theta\bigl(c(S\setminus T_\theta)-c(T_\theta\setminus S)\bigr)=\theta\bigl(c(S)-c(T_\theta)\bigr)\ge0 .
\]
(ii) Let $x^*$ be the greedy solution and $\theta$ the price per byte of the last
(partial) item. Then $x^*_i=1$ if $v_i<\theta$, $x^*_i=0$ if $v_i>\theta$, and
$\sum c_ix^*_i=\Delta C$. For any feasible $x$,
\[
  \sum_iw_i(x_i-x^*_i)=\sum_i(w_i-\theta c_i)(x_i-x^*_i)+\theta\sum_ic_i(x_i-x^*_i) .
\]
Each term of the first sum is nonnegative (if $v_i<\theta$ both factors are
$\le0$, if $v_i>\theta$ both are $\ge0$, and if $v_i=\theta$ the term is $0$). The
second sum is $\theta(\sum c_ix_i-\Delta C)\ge0$.
\end{proof}

In the integer version, price-per-byte order can be arbitrarily bad. With target
$\Delta C=2$ and three items $(c,w)=(1,1),(L,L),(2,3)$, price-per-byte order (ties
broken in favour of the earlier item) takes the first two items at cost $1+L$,
while the third item alone costs $3$. The ratio $(1+L)/3$ diverges as
$L\to\infty$. The standard repair is the \emph{guarded greedy} algorithm, which
guesses the most expensive evicted item and fills the rest in price-per-byte
order; it is known to be within a factor $2$ of the optimum for all $w_i$
\cite{csirik1991}.

Real systems evict by a key fixed in advance rather than by a price: shortest
context first (shortest-first, SF), least recently used (LRU), time-to-live
expiry (TTL), and so on. Such keys do not look at the resume probability $p_i$.

\begin{proposition}[Price-blind keys]\label{prop:blind}
Consider the covering knapsack problem with $w_i=p_i\Phi_i$.
\begin{enumerate}[label=(\roman*)]
\item Consider a rule that chooses one of two programs from any information except
  their resume probabilities. For every $R$ there are resume probabilities
  $p\in(0,1]$ under which either program alone meets the target and the program
  the rule chooses costs more than $R$ times the other. Hence no such rule has a
  constant approximation ratio.
\item If every program resumes ($p_i=1$) and the price is only the quadratic term
  of the recompute work, so that $w_i=c_i^2$, then SF equals price-per-byte order,
  costs at most $2$ times any feasible solution, and the constant $2$ cannot be
  improved.
\end{enumerate}
\end{proposition}

\begin{proof}
(i) Take two programs of the same size $c>0$ that agree in everything the rule
observes, and a target $\Delta C=c$. The rule's input does not depend on the
resume probabilities, so it drops a fixed one, $e$; call the other $o$. Set
$\Phi_e=\Phi_o=\Phi>0$, $p_e=1$ and $p_o=1/(|R|+2)$. Then
$w_e=\Phi>R\,\Phi/(|R|+2)=R\,w_o$. The optimum is at most $w_o$, so the rule's
ratio exceeds $R$.

(ii) Since $w_i/c_i=c_i$, price-per-byte order is SF. SF takes programs from the
shortest until the target is met, so it returns a feasible solution whenever the
pool holds the target. Suppose SF takes $t_1\le\dots\le t_m$ in sorted order, and
let $\ell=t_m$ and $\sigma=t_1+\dots+t_{m-1}<\Delta C$ (it did not stop earlier).
Let $S$ be any feasible solution.

\emph{Step 1: $\ell^2\le w(S)$.} Every program with $c_i<\ell$ is among the first
$m-1$, and these free only $\sigma<\Delta C$. So $S$ contains a program with
$c_i\ge\ell$.

\emph{Step 2.} Put $f(x)=x(\ell-x)$ for $x<\ell$ and $f(x)=0$ for $x\ge\ell$. Then
$f\ge0$ and $x^2\ge\ell x-f(x)$ for all $x\ge0$. Summing over $S$ and using
$c(S)\ge\Delta C$, and noting that the terms with $f(c_i)>0$ occur only among
$t_1,\dots,t_{m-1}$,
\[
  w(S)\ge\ell\,\Delta C-\sum_{j<m}t_j(\ell-t_j)=\ell(\Delta C-\sigma)+\sum_{j<m}t_j^2 .
\]
\emph{Step 3.} Since $\Delta C-\sigma>0$,
$w(\mathrm{SF})=\sum_{j<m}t_j^2+\ell^2\le w(S)-\ell(\Delta C-\sigma)+\ell^2\le w(S)+\ell^2\le2\,w(S)$.

Optimality of the constant 2: with sizes $K,K+1$ and $\Delta C=K+1$, SF takes $K$
and then $K+1$ at cost $K^2+(K+1)^2$, while $\{K+1\}$ costs $(K+1)^2$. The ratio
tends to $2$ as $K\to\infty$.
\end{proof}

\begin{example}\label{L4:ex:sf}
Even in the situation of (ii), SF is not optimal. With contexts $4,5,6$ and target
$6$, SF drops $\{4,5\}$ at cost $16+25=41$, while $\{6\}$ is feasible at cost $36$.
\end{example}

Part (i) says that when resume probabilities differ between programs, no fixed
key can stand in for the price. How large such differences are in practice is a
question about the workload and must be measured.

\section{Summary}
\begin{itemize}
\item In a closed system the arrival theorem replaces PASTA, and mean value analysis gives the mean queue; for exponential work the open formula overstates the wait.
\item Campbell's formula gives the mean resident KV, above the naive product when long sessions have long contexts.
\item Eviction is a covering knapsack: the threshold rule solves its fractional version, and a key blind to resume probabilities has no constant ratio.
\end{itemize}

\section{Exercises}

\begin{exercise}\label{L4:ex:limit}
Prove the limit in \cref{L4:rem:arrival-limit}.
\emph{Hint:} with $Z=N/\lambda$, $\pi_N(n)\propto\frac{N!}{(N-n)!\,N^n}\rho^n$, and
$\frac{N!}{(N-n)!\,N^n}=\prod_{k<n}(1-k/N)$ lies in $[0,1]$ and tends to $1$; use
dominated convergence with $\rho^n$.
\end{exercise}

\begin{exercise}
For $N=2$ and $\gamma=\mu Z$, compute $\sum_nn\,\pi_2(n)$ directly from
\eqref{L4:eq:piN} and check that it equals $Q_2$ of \cref{thm:mva}.
\emph{Hint:} $G_2=1+2/\gamma+2/\gamma^2$.
\end{exercise}

\begin{exercise}
Show that the stationary mean number of customers in M/G/$\infty$ is $\lambda\E[S]$
without using \cref{L4:prop:mginf}.
\emph{Hint:} at a delay station $W=\E[S]$, so use Little's law.
\end{exercise}

\begin{exercise}\label{L4:ex:growing}
Sessions arrive at rate $\Lambda$ and live for $T_s$, with $\kappa$ KV bytes per token. A session's context grows
linearly, $K(u)=k_0+gu$ at age $u$. Use \cref{cor:resident-kv} to show that the mean
resident KV is $\kappa\Lambda\bigl(k_0\E[T_s]+g\E[T_s^2]/2\bigr)$. For exponential $T_s$,
compare it with the naive estimate $\kappa\Lambda\E[T_s]\,(k_0+g\E[T_s]/2)$, which uses
the mean lifetime alone.
\emph{Hint:} $\E[T_s^2]=2\E[T_s]^2$ for an exponential lifetime.
\end{exercise}
